\reportpage{06 / EXPLORE, REPRODUCE AND CRITIQUE}{A guide for readers and prospective mentors}
\begin{minipage}[t]{.485\linewidth}\vspace{0pt}\setlength{\parskip}{3pt}
\subsection*{Five ways to explore}
\begin{enumerate}[leftmargin=1.4em,itemsep=5pt]
\item \textbf{Inspect Galactic evidence.} Separate the 199 VLBI entries [12] from the ESA/Gaia reconstruction. \emph{Show warped disk} selects labelled $5\times$ height; restore $1\times$ for true proportions.
\item \textbf{Resolve 100 light-years.} Open the journey preview: this takes about 128,608 model years. The recording control captures the local scene while the tab remains visible.
\item \textbf{Change the frame.} In \emph{Sun \& all planets}, compare full Galactic and 98\% co-moving paths. \emph{Inspect a planet} follows a selected body at physical radius; \emph{Whole system} restores context.
\item \textbf{Read physical quantities.} \emph{Space \& Time} shows the whole system on a years clock; selecting Earth opens a local clock experiment. The 1PN lab integrates coupled dynamics separately.
\item \textbf{Test gas separately.} The gas lab solves 2D isothermal Euler flow on a $64\times64$ grid. Height graphs density, not warp. It is not a fitted 3D/MHD Galaxy simulation.
\end{enumerate}
\end{minipage}\hfill
\begin{minipage}[t]{.485\linewidth}\vspace{0pt}\setlength{\parskip}{3pt}
\subsection*{Questions for the community}
\textbf{A feasible next test.} Which known systems have planetary orbital planes and relative velocities well enough constrained to test entry into a transit band?

\textbf{A realistic observing programme.} Which apparent gains survive exposure-time budgets, repeated visits, binaries and instrument noise?

\textbf{A second observer.} What physically feasible probe orbit would add a useful baseline beyond Earth's orbit and long-term astrometry?

\textbf{Relevant prior work.} Which studies best connect changing sightlines, nearby encounters and exoplanet detection opportunities?

\textit{I am an independent researcher without a formal supervisor. I built this software out of curiosity and would welcome expert guidance on what is useful, what is already established, and what needs correction.}
\end{minipage}
\vspace{7pt}
\subsection*{Selected sources and provenance}
\begingroup\sffamily\fontsize{8}{9.6}\selectfont
\begin{multicols}{2}
\reference{1}{ESA/Gaia/DPAC; S. Payne-Wardenaar (2025). Milky Way reconstruction. CC BY-SA 3.0 IGO; annotated here.}{https://www.esa.int/ESA_Multimedia/Images/2025/01/The_best_Milky_Way_map_by_Gaia}
\reference{2}{Finding Earth 2.0. Saved analysis, 28 Aug 2026. Original habitable-zone plot; repository and methods.}{https://github.com/Biswajit1999/finding-earth-2/tree/14828643bdd9342eeb2aff65c5dd32eef31dd2c3}
\reference{3}{Kaltenegger \& Faherty (2021). Past, present and future stars that can see Earth as a transiting exoplanet.}{https://doi.org/10.1038/s41586-021-03596-y}
\reference{4}{Gaia DR3 (2022). Catalogue release; sample selection and covariance treatment in the full methods.}{https://doi.org/10.5270/esa-qa4lep3}
\reference{5}{Reyl\'e et al. (2021). The 10 parsec sample in the Gaia era. Historical local-census normalization.}{https://doi.org/10.1051/0004-6361/202140985}
\reference{6}{Juri\'c et al. (2008). Milky Way tomography with SDSS. Basis of the adopted smooth-density scenario.}{https://arxiv.org/abs/astro-ph/0510520}
\reference{7}{Kervella, Th\'evenin \& Lovis (2017). Proxima / Alpha Centauri measurement inputs.}{https://doi.org/10.1051/0004-6361/201629930}
\reference{8}{Poggio et al. (2025). The great wave: outer-disk warp and corrugation. Not fitted by this simulation.}{https://doi.org/10.1051/0004-6361/202451668}
\reference{9}{NASA. Uncovering our Solar System's shape. Heliosphere structure and observational context.}{https://www.nasa.gov/solar-system/uncovering-our-solar-systems-shape/}
\reference{10}{Ashby (2003). Relativity in the Global Positioning System.}{https://doi.org/10.12942/lrr-2003-1}
\reference{11}{Park et al. (2021). JPL DE440/DE441 planetary and lunar ephemerides.}{https://doi.org/10.3847/1538-3881/abd414}
\reference{12}{Reid et al. (2019). Trigonometric parallaxes and Galactic spiral structure. VLBI data and arm fits.}{https://arxiv.org/abs/1910.03357}
\end{multicols}
\endgroup
\vspace{2pt}
\source{Figure 2 is the original saved plot; figures 3--5 are vector plots of project calculations. The original discussion note is preserved unchanged. Planet maps in the simulation: Solar System Scope / INOVE, CC BY 4.0, with illustrative surface and orientation choices. No institutional endorsement is implied.}

\takeaway{Open the project}{\href{https://biswajit1999.github.io/moving-observer/}{Simulation and full methods}\quad\textbar\quad\href{https://github.com/Biswajit1999/moving-observer}{Source, data and LaTeX}\quad\textbar\quad\href{https://github.com/Biswajit1999/moving-observer/issues/new/choose}{Questions and scientific criticism}}
